
\documentclass{article}
\usepackage{colortbl}
\usepackage{makecell}
\usepackage{multirow}
\usepackage{supertabular}

\begin{document}

\newcounter{utterance}

\centering \large Interaction Transcript for game `matchit\_ascii', experiment `similar\_grid\_2', episode 1 with claude{-}opus{-}4{-}8{-}azure.
\vspace{24pt}

{ \footnotesize  \setcounter{utterance}{1}
\setlength{\tabcolsep}{0pt}
\begin{supertabular}{c@{$\;$}|p{.15\linewidth}@{}p{.15\linewidth}p{.15\linewidth}p{.15\linewidth}p{.15\linewidth}p{.15\linewidth}}
    \# & $\;$A & \multicolumn{4}{c}{Game Master} & $\;\:$B\\
    \hline

    \theutterance \stepcounter{utterance}  
    & & \multicolumn{4}{p{0.6\linewidth}}{
        \cellcolor[rgb]{0.9,0.9,0.9}{
            \makecell[{{p{\linewidth}}}]{
                \texttt{\tiny{[P1$\langle$GM]}}
                \texttt{Sudjeluješ u igri pogađanja u suradnji.} \\
\texttt{Cilj je otkriti jesu li ova mreža i druga slika koju samo ja mogu vidjeti iste.} \\
\texttt{Molim te, prvo opiši svoju mrežu. Zatim ću ja dati svoj opis i možemo si međusobno postavljati pitanja o mrežama kako bismo utvrdili jesu li iste. Sada započni svoj kratki opis mreže s "OPIS:" nakon čega slijedi opis. Ne dodaj ništa drugo. Ovo je tvoja mreža:} \\
\texttt{▢ ▢ X ▢ ▢} \\
\texttt{▢ ▢ X ▢ ▢} \\
\texttt{▢ ▢ X ▢ ▢} \\
\texttt{▢ ▢ X ▢ ▢} \\
\texttt{X X X X X} \\
            }
        }
    }
    & \\ \\

    \theutterance \stepcounter{utterance}  
    & \multicolumn{4}{p{0.6\linewidth}}{
        \cellcolor[rgb]{0.9,0.9,0.9}{
            \makecell[{{p{\linewidth}}}]{
                \texttt{\tiny{[P1$\rangle$GM]}}
                \texttt{OPIS: Mreža je 5x5. U trećem stupcu (sredina) svi su kvadratići označeni X{-}om, od vrha do dna. Osim toga, cijeli donji red (peti red) je označen X{-}om. Ostali kvadratići su prazni. Oblik podsjeća na slovo "T" okrenuto naopako ili na uspravni križ koji stoji na podlozi.} \\
            }
        }
    }
    & & \\ \\

    \theutterance \stepcounter{utterance}  
    & & & \multicolumn{2}{p{0.3\linewidth}}{
        \cellcolor[rgb]{0.9,0.9,0.9}{
            \makecell[{{p{\linewidth}}}]{
                \texttt{\tiny{[GM$|$GM]}}
                \texttt{continue} \\
            }
        }
    }
    & & \\ \\

    \theutterance \stepcounter{utterance}  
    & & \multicolumn{4}{p{0.6\linewidth}}{
        \cellcolor[rgb]{0.9,0.9,0.9}{
            \makecell[{{p{\linewidth}}}]{
                \texttt{\tiny{[GM$\rangle$P2]}}
                \texttt{Sudjeluješ u igri pogađanja u suradnji.} \\
\texttt{Cilj je otkriti jesu li ova mreža i druga slika koju samo ja mogu vidjeti iste.} \\
\texttt{Molim te, prvo opiši svoju mrežu. Zatim ću ja dati svoj opis i možemo si međusobno postavljati pitanja o mrežama kako bismo utvrdili jesu li iste. Sada započni svoj kratki opis mreže s "OPIS:" nakon čega slijedi opis. Ne dodaj ništa drugo. Ovo je tvoja mreža:} \\
\texttt{▢ ▢ X ▢ ▢} \\
\texttt{▢ ▢ X ▢ ▢} \\
\texttt{▢ ▢ ▢ ▢ ▢} \\
\texttt{▢ ▢ X X ▢} \\
\texttt{X X X X X} \\
            }
        }
    }
    & \\ \\

    \theutterance \stepcounter{utterance}  
    & & & \multicolumn{4}{p{0.6\linewidth}}{
        \cellcolor[rgb]{0.9,0.9,0.9}{
            \makecell[{{p{\linewidth}}}]{
                \texttt{\tiny{[GM$\langle$P2]}}
                \texttt{OPIS: Imam mrežu 5x5. U prvom redu je X na trećem stupcu. U drugom redu je X na trećem stupcu. Treći red je prazan. U četvrtom redu su X na trećem i četvrtom stupcu. U petom (zadnjem) redu su svi X (svih pet polja).} \\
            }
        }
    }
     \\ \\

    \theutterance \stepcounter{utterance}  
    & & & \multicolumn{2}{p{0.3\linewidth}}{
        \cellcolor[rgb]{0.9,0.9,0.9}{
            \makecell[{{p{\linewidth}}}]{
                \texttt{\tiny{[GM$|$GM]}}
                \texttt{continue} \\
            }
        }
    }
    & & \\ \\

    \theutterance \stepcounter{utterance}  
    & & \multicolumn{4}{p{0.6\linewidth}}{
        \cellcolor[rgb]{0.9,0.9,0.9}{
            \makecell[{{p{\linewidth}}}]{
                \texttt{\tiny{[GM$\rangle$P2]}}
                \texttt{Ovo je mojOPIS: Mreža je 5x5. U trećem stupcu (sredina) svi su kvadratići označeni X{-}om, od vrha do dna. Osim toga, cijeli donji red (peti red) je označen X{-}om. Ostali kvadratići su prazni. Oblik podsjeća na slovo "T" okrenuto naopako ili na uspravni križ koji stoji na podlozi.} \\
\texttt{Sada postavi pitanje kako bi saznao nove aspekte moje mreže koji se možda razlikuju od tvoje mreže. Započni s "PITANJE:" i ne dodavaj ništa drugo.} \\
            }
        }
    }
    & \\ \\

    \theutterance \stepcounter{utterance}  
    & & & \multicolumn{4}{p{0.6\linewidth}}{
        \cellcolor[rgb]{0.9,0.9,0.9}{
            \makecell[{{p{\linewidth}}}]{
                \texttt{\tiny{[GM$\langle$P2]}}
                \texttt{PITANJE: U tvojoj mreži, ima li X na četvrtom stupcu u četvrtom redu (dakle desno od sredine u pretposljednjem redu)?} \\
            }
        }
    }
     \\ \\

    \theutterance \stepcounter{utterance}  
    & & & \multicolumn{2}{p{0.3\linewidth}}{
        \cellcolor[rgb]{0.9,0.9,0.9}{
            \makecell[{{p{\linewidth}}}]{
                \texttt{\tiny{[GM$|$GM]}}
                \texttt{continue} \\
            }
        }
    }
    & & \\ \\

    \theutterance \stepcounter{utterance}  
    & & \multicolumn{4}{p{0.6\linewidth}}{
        \cellcolor[rgb]{0.9,0.9,0.9}{
            \makecell[{{p{\linewidth}}}]{
                \texttt{\tiny{[P1$\langle$GM]}}
                \texttt{Ovo je mojOPIS: Imam mrežu 5x5. U prvom redu je X na trećem stupcu. U drugom redu je X na trećem stupcu. Treći red je prazan. U četvrtom redu su X na trećem i četvrtom stupcu. U petom (zadnjem) redu su svi X (svih pet polja).} \\
\texttt{PITANJE: U tvojoj mreži, ima li X na četvrtom stupcu u četvrtom redu (dakle desno od sredine u pretposljednjem redu)?Započni svoj odgovor s "ODGOVOR:" i nemoj dodati ništa drugo.} \\
            }
        }
    }
    & \\ \\

    \theutterance \stepcounter{utterance}  
    & \multicolumn{4}{p{0.6\linewidth}}{
        \cellcolor[rgb]{0.9,0.9,0.9}{
            \makecell[{{p{\linewidth}}}]{
                \texttt{\tiny{[P1$\rangle$GM]}}
                \texttt{ODGOVOR: Ne, u mojoj mreži četvrti red ima X samo na trećem stupcu, a četvrti stupac u četvrtom redu je prazan. Također, u mojoj mreži treći red NIJE prazan — ima X na trećem stupcu. Čini se da naše mreže nisu iste.} \\
            }
        }
    }
    & & \\ \\

    \theutterance \stepcounter{utterance}  
    & & & \multicolumn{2}{p{0.3\linewidth}}{
        \cellcolor[rgb]{0.9,0.9,0.9}{
            \makecell[{{p{\linewidth}}}]{
                \texttt{\tiny{[GM$|$GM]}}
                \texttt{continue} \\
            }
        }
    }
    & & \\ \\

    \theutterance \stepcounter{utterance}  
    & & \multicolumn{4}{p{0.6\linewidth}}{
        \cellcolor[rgb]{0.9,0.9,0.9}{
            \makecell[{{p{\linewidth}}}]{
                \texttt{\tiny{[P1$\langle$GM]}}
                \texttt{Sada postavi pitanje kako bi saznao nove aspekte moje mreže koji se možda razlikuju od tvoje mreže. Započni s "PITANJE:" i ne dodavaj ništa drugo.} \\
            }
        }
    }
    & \\ \\

    \theutterance \stepcounter{utterance}  
    & \multicolumn{4}{p{0.6\linewidth}}{
        \cellcolor[rgb]{0.9,0.9,0.9}{
            \makecell[{{p{\linewidth}}}]{
                \texttt{\tiny{[P1$\rangle$GM]}}
                \texttt{PITANJE: U tvojoj mreži, ima li X na trećem stupcu u trećem redu (dakle u samom središtu mreže)?} \\
            }
        }
    }
    & & \\ \\

    \theutterance \stepcounter{utterance}  
    & & & \multicolumn{2}{p{0.3\linewidth}}{
        \cellcolor[rgb]{0.9,0.9,0.9}{
            \makecell[{{p{\linewidth}}}]{
                \texttt{\tiny{[GM$|$GM]}}
                \texttt{continue} \\
            }
        }
    }
    & & \\ \\

    \theutterance \stepcounter{utterance}  
    & & \multicolumn{4}{p{0.6\linewidth}}{
        \cellcolor[rgb]{0.9,0.9,0.9}{
            \makecell[{{p{\linewidth}}}]{
                \texttt{\tiny{[GM$\rangle$P2]}}
                \texttt{ODGOVOR: Ne, u mojoj mreži četvrti red ima X samo na trećem stupcu, a četvrti stupac u četvrtom redu je prazan. Također, u mojoj mreži treći red NIJE prazan — ima X na trećem stupcu. Čini se da naše mreže nisu iste.} \\
\texttt{PITANJE: U tvojoj mreži, ima li X na trećem stupcu u trećem redu (dakle u samom središtu mreže)?Započni svoj odgovor s "ODGOVOR:" i nemoj dodati ništa drugo.} \\
            }
        }
    }
    & \\ \\

    \theutterance \stepcounter{utterance}  
    & & & \multicolumn{4}{p{0.6\linewidth}}{
        \cellcolor[rgb]{0.9,0.9,0.9}{
            \makecell[{{p{\linewidth}}}]{
                \texttt{\tiny{[GM$\langle$P2]}}
                \texttt{ODGOVOR: Ne, u mojoj mreži treći red je potpuno prazan, uključujući i središte (treći stupac, treći red). Osim toga, moja mreža ima X na četvrtom stupcu u četvrtom redu, dok tvoja nema. Naše mreže se razlikuju.} \\
            }
        }
    }
     \\ \\

    \theutterance \stepcounter{utterance}  
    & & & \multicolumn{2}{p{0.3\linewidth}}{
        \cellcolor[rgb]{0.9,0.9,0.9}{
            \makecell[{{p{\linewidth}}}]{
                \texttt{\tiny{[GM$|$GM]}}
                \texttt{continue} \\
            }
        }
    }
    & & \\ \\

    \theutterance \stepcounter{utterance}  
    & & \multicolumn{4}{p{0.6\linewidth}}{
        \cellcolor[rgb]{0.9,0.9,0.9}{
            \makecell[{{p{\linewidth}}}]{
                \texttt{\tiny{[GM$\rangle$P2]}}
                \texttt{Sada postavi pitanje kako bi saznao nove aspekte moje mreže koji se možda razlikuju od tvoje mreže. Započni s "PITANJE:" i ne dodavaj ništa drugo.} \\
            }
        }
    }
    & \\ \\

    \theutterance \stepcounter{utterance}  
    & & & \multicolumn{4}{p{0.6\linewidth}}{
        \cellcolor[rgb]{0.9,0.9,0.9}{
            \makecell[{{p{\linewidth}}}]{
                \texttt{\tiny{[GM$\langle$P2]}}
                \texttt{PITANJE: Potvrđuješ li da u tvojoj mreži nema nijednog X izvan trećeg stupca i zadnjeg reda (dakle da su svi ostali kvadratići prazni)?} \\
            }
        }
    }
     \\ \\

    \theutterance \stepcounter{utterance}  
    & & & \multicolumn{2}{p{0.3\linewidth}}{
        \cellcolor[rgb]{0.9,0.9,0.9}{
            \makecell[{{p{\linewidth}}}]{
                \texttt{\tiny{[GM$|$GM]}}
                \texttt{continue} \\
            }
        }
    }
    & & \\ \\

    \theutterance \stepcounter{utterance}  
    & & \multicolumn{4}{p{0.6\linewidth}}{
        \cellcolor[rgb]{0.9,0.9,0.9}{
            \makecell[{{p{\linewidth}}}]{
                \texttt{\tiny{[P1$\langle$GM]}}
                \texttt{ODGOVOR: Ne, u mojoj mreži treći red je potpuno prazan, uključujući i središte (treći stupac, treći red). Osim toga, moja mreža ima X na četvrtom stupcu u četvrtom redu, dok tvoja nema. Naše mreže se razlikuju.} \\
\texttt{PITANJE: Potvrđuješ li da u tvojoj mreži nema nijednog X izvan trećeg stupca i zadnjeg reda (dakle da su svi ostali kvadratići prazni)?Započni svoj odgovor s "ODGOVOR:" i nemoj dodati ništa drugo.} \\
            }
        }
    }
    & \\ \\

    \theutterance \stepcounter{utterance}  
    & \multicolumn{4}{p{0.6\linewidth}}{
        \cellcolor[rgb]{0.9,0.9,0.9}{
            \makecell[{{p{\linewidth}}}]{
                \texttt{\tiny{[P1$\rangle$GM]}}
                \texttt{ODGOVOR: Da, potvrđujem. U mojoj mreži svi X{-}ovi su isključivo u trećem stupcu (od vrha do dna) i u zadnjem redu. Svi ostali kvadratići su prazni. Budući da tvoja mreža ima prazan treći red i X na četvrtom stupcu u četvrtom redu, naše mreže se definitivno razlikuju.} \\
            }
        }
    }
    & & \\ \\

    \theutterance \stepcounter{utterance}  
    & & & \multicolumn{2}{p{0.3\linewidth}}{
        \cellcolor[rgb]{0.9,0.9,0.9}{
            \makecell[{{p{\linewidth}}}]{
                \texttt{\tiny{[GM$|$GM]}}
                \texttt{continue} \\
            }
        }
    }
    & & \\ \\

    \theutterance \stepcounter{utterance}  
    & & \multicolumn{4}{p{0.6\linewidth}}{
        \cellcolor[rgb]{0.9,0.9,0.9}{
            \makecell[{{p{\linewidth}}}]{
                \texttt{\tiny{[P1$\langle$GM]}}
                \texttt{Sada postavi pitanje kako bi saznao nove aspekte moje mreže koji se možda razlikuju od tvoje mreže. Započni s "PITANJE:" i ne dodavaj ništa drugo.} \\
            }
        }
    }
    & \\ \\

    \theutterance \stepcounter{utterance}  
    & \multicolumn{4}{p{0.6\linewidth}}{
        \cellcolor[rgb]{0.9,0.9,0.9}{
            \makecell[{{p{\linewidth}}}]{
                \texttt{\tiny{[P1$\rangle$GM]}}
                \texttt{PITANJE: U tvojoj mreži, u zadnjem redu, jesu li svih pet polja X, uključujući i prvo i peto polje (dakle lijevi i desni rub)?} \\
            }
        }
    }
    & & \\ \\

    \theutterance \stepcounter{utterance}  
    & & & \multicolumn{2}{p{0.3\linewidth}}{
        \cellcolor[rgb]{0.9,0.9,0.9}{
            \makecell[{{p{\linewidth}}}]{
                \texttt{\tiny{[GM$|$GM]}}
                \texttt{continue} \\
            }
        }
    }
    & & \\ \\

    \theutterance \stepcounter{utterance}  
    & & \multicolumn{4}{p{0.6\linewidth}}{
        \cellcolor[rgb]{0.9,0.9,0.9}{
            \makecell[{{p{\linewidth}}}]{
                \texttt{\tiny{[GM$\rangle$P2]}}
                \texttt{ODGOVOR: Da, potvrđujem. U mojoj mreži svi X{-}ovi su isključivo u trećem stupcu (od vrha do dna) i u zadnjem redu. Svi ostali kvadratići su prazni. Budući da tvoja mreža ima prazan treći red i X na četvrtom stupcu u četvrtom redu, naše mreže se definitivno razlikuju.} \\
\texttt{PITANJE: U tvojoj mreži, u zadnjem redu, jesu li svih pet polja X, uključujući i prvo i peto polje (dakle lijevi i desni rub)?Započni svoj odgovor s "ODGOVOR:" i nemoj dodati ništa drugo.} \\
            }
        }
    }
    & \\ \\

    \theutterance \stepcounter{utterance}  
    & & & \multicolumn{4}{p{0.6\linewidth}}{
        \cellcolor[rgb]{0.9,0.9,0.9}{
            \makecell[{{p{\linewidth}}}]{
                \texttt{\tiny{[GM$\langle$P2]}}
                \texttt{ODGOVOR: Da, u mojoj mreži je zadnji red potpuno popunjen — svih pet polja su X, uključujući i lijevi i desni rub. To se poklapa s tvojom mrežom, no ostale razlike (prazan treći red i X na četvrtom stupcu u četvrtom redu kod mene) potvrđuju da naše mreže nisu iste.} \\
            }
        }
    }
     \\ \\

    \theutterance \stepcounter{utterance}  
    & & & \multicolumn{2}{p{0.3\linewidth}}{
        \cellcolor[rgb]{0.9,0.9,0.9}{
            \makecell[{{p{\linewidth}}}]{
                \texttt{\tiny{[GM$|$GM]}}
                \texttt{continue} \\
            }
        }
    }
    & & \\ \\

    \theutterance \stepcounter{utterance}  
    & & \multicolumn{4}{p{0.6\linewidth}}{
        \cellcolor[rgb]{0.9,0.9,0.9}{
            \makecell[{{p{\linewidth}}}]{
                \texttt{\tiny{[GM$\rangle$P2]}}
                \texttt{Sada postavi pitanje kako bi saznao nove aspekte moje mreže koji se možda razlikuju od tvoje mreže. Započni s "PITANJE:" i ne dodavaj ništa drugo.} \\
            }
        }
    }
    & \\ \\

    \theutterance \stepcounter{utterance}  
    & & & \multicolumn{4}{p{0.6\linewidth}}{
        \cellcolor[rgb]{0.9,0.9,0.9}{
            \makecell[{{p{\linewidth}}}]{
                \texttt{\tiny{[GM$\langle$P2]}}
                \texttt{PITANJE: U tvojoj mreži, u prvom i drugom redu, ima li X samo na trećem stupcu (dakle poklapa li se to s mojom mrežom u ta dva reda)?} \\
            }
        }
    }
     \\ \\

    \theutterance \stepcounter{utterance}  
    & & & \multicolumn{2}{p{0.3\linewidth}}{
        \cellcolor[rgb]{0.9,0.9,0.9}{
            \makecell[{{p{\linewidth}}}]{
                \texttt{\tiny{[GM$|$GM]}}
                \texttt{continue} \\
            }
        }
    }
    & & \\ \\

    \theutterance \stepcounter{utterance}  
    & & \multicolumn{4}{p{0.6\linewidth}}{
        \cellcolor[rgb]{0.9,0.9,0.9}{
            \makecell[{{p{\linewidth}}}]{
                \texttt{\tiny{[P1$\langle$GM]}}
                \texttt{ODGOVOR: Da, u mojoj mreži je zadnji red potpuno popunjen — svih pet polja su X, uključujući i lijevi i desni rub. To se poklapa s tvojom mrežom, no ostale razlike (prazan treći red i X na četvrtom stupcu u četvrtom redu kod mene) potvrđuju da naše mreže nisu iste.} \\
\texttt{PITANJE: U tvojoj mreži, u prvom i drugom redu, ima li X samo na trećem stupcu (dakle poklapa li se to s mojom mrežom u ta dva reda)?Započni svoj odgovor s "ODGOVOR:" i nemoj dodati ništa drugo.} \\
            }
        }
    }
    & \\ \\

    \theutterance \stepcounter{utterance}  
    & \multicolumn{4}{p{0.6\linewidth}}{
        \cellcolor[rgb]{0.9,0.9,0.9}{
            \makecell[{{p{\linewidth}}}]{
                \texttt{\tiny{[P1$\rangle$GM]}}
                \texttt{ODGOVOR: Da, u mojoj mreži prvi i drugi red imaju X isključivo na trećem stupcu, što se poklapa s tvojom mrežom u ta dva reda. Međutim, zbog razlika u trećem i četvrtom redu, naše mreže ukupno nisu iste.} \\
            }
        }
    }
    & & \\ \\

    \theutterance \stepcounter{utterance}  
    & & & \multicolumn{2}{p{0.3\linewidth}}{
        \cellcolor[rgb]{0.9,0.9,0.9}{
            \makecell[{{p{\linewidth}}}]{
                \texttt{\tiny{[GM$|$GM]}}
                \texttt{continue} \\
            }
        }
    }
    & & \\ \\

    \theutterance \stepcounter{utterance}  
    & & \multicolumn{4}{p{0.6\linewidth}}{
        \cellcolor[rgb]{0.9,0.9,0.9}{
            \makecell[{{p{\linewidth}}}]{
                \texttt{\tiny{[P1$\langle$GM]}}
                \texttt{Sada postavi pitanje kako bi saznao nove aspekte moje mreže koji se možda razlikuju od tvoje mreže. Započni s "PITANJE:" i ne dodavaj ništa drugo.} \\
            }
        }
    }
    & \\ \\

    \theutterance \stepcounter{utterance}  
    & \multicolumn{4}{p{0.6\linewidth}}{
        \cellcolor[rgb]{0.9,0.9,0.9}{
            \makecell[{{p{\linewidth}}}]{
                \texttt{\tiny{[P1$\rangle$GM]}}
                \texttt{PITANJE: U tvojoj mreži, u četvrtom redu, potvrđuješ li da su X samo na trećem i četvrtom stupcu, a da su prvi, drugi i peti stupac prazni?} \\
            }
        }
    }
    & & \\ \\

    \theutterance \stepcounter{utterance}  
    & & & \multicolumn{2}{p{0.3\linewidth}}{
        \cellcolor[rgb]{0.9,0.9,0.9}{
            \makecell[{{p{\linewidth}}}]{
                \texttt{\tiny{[GM$|$GM]}}
                \texttt{continue} \\
            }
        }
    }
    & & \\ \\

    \theutterance \stepcounter{utterance}  
    & & \multicolumn{4}{p{0.6\linewidth}}{
        \cellcolor[rgb]{0.9,0.9,0.9}{
            \makecell[{{p{\linewidth}}}]{
                \texttt{\tiny{[GM$\rangle$P2]}}
                \texttt{ODGOVOR: Da, u mojoj mreži prvi i drugi red imaju X isključivo na trećem stupcu, što se poklapa s tvojom mrežom u ta dva reda. Međutim, zbog razlika u trećem i četvrtom redu, naše mreže ukupno nisu iste.} \\
\texttt{PITANJE: U tvojoj mreži, u četvrtom redu, potvrđuješ li da su X samo na trećem i četvrtom stupcu, a da su prvi, drugi i peti stupac prazni?Započni svoj odgovor s "ODGOVOR:" i nemoj dodati ništa drugo.} \\
            }
        }
    }
    & \\ \\

    \theutterance \stepcounter{utterance}  
    & & & \multicolumn{4}{p{0.6\linewidth}}{
        \cellcolor[rgb]{0.9,0.9,0.9}{
            \makecell[{{p{\linewidth}}}]{
                \texttt{\tiny{[GM$\langle$P2]}}
                \texttt{ODGOVOR: Da, u mojoj mreži četvrti red ima X isključivo na trećem i četvrtom stupcu, dok su prvi, drugi i peti stupac prazni. Budući da tvoja mreža u četvrtom redu ima X samo na trećem stupcu, i budući da je kod mene treći red prazan a kod tebe nije, naše mreže se razlikuju.} \\
            }
        }
    }
     \\ \\

    \theutterance \stepcounter{utterance}  
    & & & \multicolumn{2}{p{0.3\linewidth}}{
        \cellcolor[rgb]{0.9,0.9,0.9}{
            \makecell[{{p{\linewidth}}}]{
                \texttt{\tiny{[GM$|$GM]}}
                \texttt{continue} \\
            }
        }
    }
    & & \\ \\

    \theutterance \stepcounter{utterance}  
    & & \multicolumn{4}{p{0.6\linewidth}}{
        \cellcolor[rgb]{0.9,0.9,0.9}{
            \makecell[{{p{\linewidth}}}]{
                \texttt{\tiny{[GM$\rangle$P2]}}
                \texttt{Sada donesi odluku. Što misliš: jesu li tvoja mreža i druga opisana mreža iste? Napiši "ODLUKA: ista mreža" ako misliš da su iste mreže ili "ODLUKA: različite mreže" ako misliš da su različite mreže. Ne dodaj ništa drugo.} \\
            }
        }
    }
    & \\ \\

    \theutterance \stepcounter{utterance}  
    & & & \multicolumn{4}{p{0.6\linewidth}}{
        \cellcolor[rgb]{0.9,0.9,0.9}{
            \makecell[{{p{\linewidth}}}]{
                \texttt{\tiny{[GM$\langle$P2]}}
                \texttt{ODLUKA: različite mreže} \\
            }
        }
    }
     \\ \\

    \theutterance \stepcounter{utterance}  
    & & & \multicolumn{2}{p{0.3\linewidth}}{
        \cellcolor[rgb]{0.9,0.9,0.9}{
            \makecell[{{p{\linewidth}}}]{
                \texttt{\tiny{[GM$|$GM]}}
                \texttt{continue} \\
            }
        }
    }
    & & \\ \\

    \theutterance \stepcounter{utterance}  
    & & & \multicolumn{2}{p{0.3\linewidth}}{
        \cellcolor[rgb]{0.9,0.9,0.9}{
            \makecell[{{p{\linewidth}}}]{
                \texttt{\tiny{[GM$|$GM]}}
                \texttt{success} \\
            }
        }
    }
    & & \\ \\

    \theutterance \stepcounter{utterance}  
    & & \multicolumn{4}{p{0.6\linewidth}}{
        \cellcolor[rgb]{0.9,0.9,0.9}{
            \makecell[{{p{\linewidth}}}]{
                \texttt{\tiny{[P1$\langle$GM]}}
                \texttt{ODGOVOR: Da, u mojoj mreži četvrti red ima X isključivo na trećem i četvrtom stupcu, dok su prvi, drugi i peti stupac prazni. Budući da tvoja mreža u četvrtom redu ima X samo na trećem stupcu, i budući da je kod mene treći red prazan a kod tebe nije, naše mreže se razlikuju.} \\
\texttt{Sada donesi odluku. Što misliš: jesu li tvoja mreža i druga opisana mreža iste? Napiši "ODLUKA: ista mreža" ako misliš da su iste mreže ili "ODLUKA: različite mreže" ako misliš da su različite mreže. Ne dodaj ništa drugo.} \\
            }
        }
    }
    & \\ \\

    \theutterance \stepcounter{utterance}  
    & \multicolumn{4}{p{0.6\linewidth}}{
        \cellcolor[rgb]{0.9,0.9,0.9}{
            \makecell[{{p{\linewidth}}}]{
                \texttt{\tiny{[P1$\rangle$GM]}}
                \texttt{ODLUKA: različite mreže} \\
            }
        }
    }
    & & \\ \\

    \theutterance \stepcounter{utterance}  
    & & & \multicolumn{2}{p{0.3\linewidth}}{
        \cellcolor[rgb]{0.9,0.9,0.9}{
            \makecell[{{p{\linewidth}}}]{
                \texttt{\tiny{[GM$|$GM]}}
                \texttt{continue} \\
            }
        }
    }
    & & \\ \\

    \theutterance \stepcounter{utterance}  
    & & & \multicolumn{2}{p{0.3\linewidth}}{
        \cellcolor[rgb]{0.9,0.9,0.9}{
            \makecell[{{p{\linewidth}}}]{
                \texttt{\tiny{[GM$|$GM]}}
                \texttt{success} \\
            }
        }
    }
    & & \\ \\

\end{supertabular}
}

\end{document}
